\scanpage{46}
\noindent\hfill atkin \deleted{5} \handwritten{\mspage{6}}

\noindent\deleted{\unsure{2|1}}\handwritten{3}.2. In 1962*O'Brien and I obtained a result which generalises

\noindent\handwritten{\rule{\linewidth}{0.4pt}}

\noindent *The results of \handwritten{\S}\deleted{\unsure{2|1}}\handwritten{3}.2 appear in ATKIN and O'BRIEN(1967). In the present context it seems natural to use real time.

\noindent\handwritten{\rule{\linewidth}{0.4pt}}

\noindent Newman's \handwritten{(3,1)}:
\[
\begin{aligned}
&\qquad\handwritten{\text{For all}\quad\alpha\geq1\quad\text{there exists a constant}\quad k_{\alpha}\quad\text{not \unsure{divisible|diisible}}}\\
&\handwritten{\text{by}\quad13\quad\text{such that for}\quad n\geq1\quad\text{we have}}\\
&\qquad\handwritten{c(13^{\alpha+1}n)\equiv k_{\alpha}\,c(13^{\alpha}n)\quad(\mathrm{mod}\;13^{\alpha}),}
\end{aligned}
\tag*{\handwritten{(3.7)}}
\]

\noindent Specifically,we had at this stage proved \handwritten{(3.7)} for $\alpha\handwritten{\leq}4$ by detailed hand computation using the irrational modular equation for $13$,and conjectured it for all $\alpha$ . It appeared to us at the time that the proof might depend on some very deep properties of certain numerical coefficients of the form $\handwritten{a_{i}\mathbin{.}13^{b_{i}}}$,and our subsequent proof of \handwritten{(\deleted{\unsure{6|1|2}}3,7)} for all $\alpha$ in 1963 arose from a fortunate numerical mistake. In fact one of us used \unsure{\textasciigrave|\textquotesingle}some incorrect values of the \uline{$a_{i}$} and still obtained \handwritten{(3,7)},suggesting that only the values of \uline{$b_{i}$} were relevant. This turned out to be the case. It thus appears that occasionally it is better \uline{not} to have a reliable computer; but this is the only instance I can quote from personal experience.

It is clear from \handwritten{(3.7)} that if we can find an ``accident'' of the form $c(13^{\alpha}n_{\unsure{0|o}})\handwritten{\equiv}0\;(\mathrm{mod}\;13^{\beta})$ for some $\beta\handwritten{\leq}\alpha$,then $c(13^{\alpha+n}n_{\unsure{0|o|i}})\handwritten{\equiv}0\;(\mathrm{mod}\;13^{\beta})$ for $n\handwritten{\geq}0$,so that \uline{$c(n)$} is divisible by $13^{\beta}$ infinitely often. After my appointment to Atlas in 1964 I computed \uline{$c(13^{6}n)$} modulo $13^{6}$ for $1\handwritten{\leq}n\handwritten{\leq}40960$ with output of the form
\diagram{$\begin{array}{ccc}
n & \qquad a(n)\qquad & b(n)
\end{array}$}
denoting that $c(13^{6}n)\handwritten{\equiv}a(n)\mathbin{.}13^{b(n)}\quad(\mathrm{mod}\;13^{6})$ ,with \uline{$a(n)$} prime to $13$. The largest $b(n)$ found was a disappointing $3$; there seemed a priori to be a reasonable chance of $4$ if the values of \uline{$c(13^{6}n)$} were random $(\mathrm{mod}\;13^{4})$ since $40960/\;13^{4}\handwritten{\bumpeq}1.4$. However an inspection of the results exposed some very interesting features. Thus $b(7)=$\ \parbox[c]{12mm}{\handwritten{\figure{\tikz[baseline=0pt,x=.28mm,y=.28mm,line cap=round]{\draw (0,0)..controls(5,6)and(9,6)..(12,2)..controls(19,-6)and(29,-6)..(37,6); \draw (1,.2)..controls(7,1.5)and(11,3)..(12,2); \draw (23,-2)..controls(30,-2)and(34,2)..(37,6);}}}}

% archive sheet number: 45
